\section*{Erd\H{o}s Problem 941}

\paragraph{Statement and result.}
Is every sufficiently large positive integer a sum of at most three powerful numbers? A positive integer is powerful if every prime dividing it has exponent at least two. The answer is yes, first proved by Heath-Brown. We reconstruct the deduction using Moroz's later quadratic-form theorem, explicitly isolating that analytic input and supplying the elementary cases, including the integers with a bounded odd part and an arbitrarily large power of four.

\paragraph{Two precisely stated external inputs.}
We use the classical three-squares theorem: a nonnegative integer is a sum of three integer squares if and only if it is not of the form \(4^j(8b+7)\), with \(j,b\geq0\).

The second input is the following consequence of Moroz's Theorem 2. For any fixed prime \(q\equiv5\pmod8\), every sufficiently large integer \(n\equiv7\pmod8\) is represented by
\[
 n=x^2+y^2+q^3z^2,\qquad x,y,z\in\mathbb Z. \tag{1}
\]
Moroz proves a positive lower bound for its representation number using estimates for ternary quadratic forms. That analytic argument is external here. We use only \(q=5\), so (1) has the fixed coefficient 125 and a single fixed threshold \(N_0\).

\paragraph{Why these forms give powerful summands.}
Every nonzero integer square is powerful. So is \(125z^2\) when \(z\ne0\): its exponent at 5 is \(3+2v_5(z)\geq3\), and every other prime exponent is even. More generally, multiplying any powerful integer by an integer square preserves powerfulness. A zero coordinate contributes no summand; it is omitted rather than called a positive powerful number.

\paragraph{The elementary cases.}
If \(n\) is not of the excluded three-squares form, the first external input already represents it by at most three positive squares, after omitting zeros. This gives the desired representation.

We also prove that every integer of the form
\[
 n=4^j m,\qquad j\geq1,\qquad m\text{ odd}, \tag{2}
\]
has the required representation, without an asymptotic threshold. Since \(v_2(2m)=1\), the integer \(2m\) is not excluded by the three-squares theorem. Write
\[
 2m=u^2+v^2+w^2.
\]
Modulo four, precisely two of \(u,v,w\) are odd. Relabel so that \(u,v\) are odd and \(w\) is even. Then
\[
 m=\left(\frac{u+v}{2}\right)^2+
   \left(\frac{u-v}{2}\right)^2+
   2\left(\frac w2\right)^2. \tag{3}
\]
All displayed quantities in parentheses are integers. Multiplying (3) by \(4^j\) produces two squares and a term of the form
\(2^{2j+1}t^2\). If this last term is nonzero, its exponent at 2 is at least three, while every odd prime has even exponent. Thus it too is powerful. This proves (2).

\paragraph{Completing the residue-class argument.}
The only integers not yet covered by the elementary cases are those with \(n\equiv7\pmod8\). Indeed, every excluded three-squares integer is \(4^j(8b+7)\); those with \(j\geq1\) are covered by (2), and the others have \(j=0\). If \(n\geq N_0\), equation (1) with \(q=5\) represents every remaining integer. Its nonzero summands are powerful by the earlier prime-exponent check.

Therefore every \(n\geq N_0\) is a sum of at most three positive powerful integers. There is no uniformity problem involving the reduced odd part: the entire \(j\geq1\) case was settled independently in (2), so the analytic theorem is applied directly to the original integer \(n\), with one fixed threshold.

\paragraph{Scope and attribution.}
The original affirmative theorem is D. R. Heath-Brown's \emph{Ternary quadratic forms and sums of three square-full numbers}, S\'eminaire de Th\'eorie des Nombres, Paris 1986--87, pp. 137--163. The precise alternate input used here is B. Z. Moroz, \emph{On the representation of large integers by integral ternary positive definite quadratic forms}, Ast\'erisque 209 (1992), Theorem 2 on p. 276 and the following corollary: \url{https://www.numdam.org/article/AST_1992__209__275_0.pdf}. See also \url{https://www.erdosproblems.com/941}, checked 24 September 2026.

This is a complete deduction relative to the stated three-squares and Moroz theorems. It supplies all residue cases and the prime-exponent checks, but does not reproduce the modular-form estimates behind (1), supply an effective numerical threshold, or claim a new result or local Lean verification.

\textbf{Completion rate estimate: 100\%.}
